\chapter{Harmonogram realizacji projektu}

Projekt można było wykonać w ciągu jednego tygodnia. Prace zostały rozdzielone na kilka prostych etapów. Najpierw wybrano temat i określono zakres funkcji. Następnie przygotowano modele danych, bazę, migracje i dane startowe. Później wykonano interfejs, obsługę operacji CRUD, walidację oraz końcowe poprawki wyglądu.

\begin{figure}[H]
\centering
\includegraphics[width=0.95\textwidth]{figures/gantt.eps}
\caption{Orientacyjny harmonogram realizacji projektu w ciągu tygodnia}
\end{figure}

\begin{table}[H]
\centering
\caption{Etapy realizacji projektu}
\begin{tabular}{|p{3cm}|p{10cm}|}
\hline
\textbf{Etap} & \textbf{Opis prac} \\
\hline
Dzień 1 & Wybór tematu, analiza wymagań i zaplanowanie tabel. \\
\hline
Dzień 2 & Utworzenie projektu WPF, modeli i kontekstu EF Core. \\
\hline
Dzień 3 & Przygotowanie migracji, repozytoriów, serwisów i danych startowych. \\
\hline
Dzień 4 & Wykonanie ViewModelu, komend i podstawowej nawigacji. \\
\hline
Dzień 5 & Dodanie operacji CRUD, płatności i dostępności pokoi. \\
\hline
Dzień 6 & Poprawa panelu głównego, modułu usług i stylizacji. \\
\hline
Dzień 7 & Testowanie, poprawki, przygotowanie dokumentacji. \\
\hline
\end{tabular}
\end{table}

Harmonogram jest orientacyjny. W praktyce część zadań mogła być wykonywana równolegle, na przykład poprawki interfejsu i poprawki walidacji danych.
